% ---
% id: receiver/rem:edgemeaning
% kind: rem
% title: "Meaning of an edge"
% summary: >-
%   An edge fixes the order in which two elements act; moving light between
%   modes is done by the elements' matrices. Components drawn on different arms
%   carry no edge and have disjoint supports, so support comparability
%   constrains only an ambiguous drawing.
% prominence: standard
% status: accepted
% uses: []
% mentions: [receiver/component-catalogue/def:catalogue-entries, prop:mode-algebra]
% assumes: []
% refs: []
% provenance: newly-derived
% lean:
% ---

\begin{remark}\label{rem:edgemeaning}
An edge $e\to f$ fixes the order in which $e$ and $f$ act: $e$ first. Moving
light from one mode to another is done by the elements' matrices, for
instance by a mode-permutation element (entry C9 of
Table~\ref{tab:catalogue}), not by the edge; between elements with disjoint
supports the order the edge fixes does not by itself change the compiled
matrix, since they commute (Proposition~\ref{prop:mode-algebra}\,(iii)); such
an edge matters only through the order it induces, by reachability, between
elements whose supports intersect. A user who
draws two components in series draws the edge; a user who draws them on
different arms of a beamsplitter draws none, and (E3) then requires their mode
supports to be disjoint, which, on different arms, they are. The condition
therefore constrains exactly the drawings that were ambiguous.
\end{remark}
